\section{Gesture-based gaming}
\subsection{Introduction}
The use of gestures as a mode of interaction has been explored in areas such as Human-Computer Interaction and Virtual Environments. However, very little work has been done to investigate their potential uses within gaming~\citep{Krat07}. Games, unlike much other software, are used for their intrinsic enjoyment value. This means that it is often easier to entice users to use novel technology or modalities of interaction within a gaming context~\citep{Star00}. While there has been growing interest in technologies that would enable such gesture-based gaming, such as Sony's EyeToy, Microsoft's Kinect and Nintendo's Wiimote, very little has been done to assess how the inclusion of gestures will affect gaming~\citep{Wils09}. This section will discuss the limited literature available on this topic.

\subsection{Previous research on gesture-based gaming}
Previous research in this area has primarily used Computer Vision or accelerometer-based means for detecting gestures~\citep{Kang04,Payn06,Wils09}. Thus the discussion below primarily pertains to spatial three-dimensional gestures. 

Studies done in this area have generally focused on how to implement gesture-based gaming rather than considering whether this interaction style increases the enjoyment and immersion a player feels during game-play~\citep{Kang04,Krat07}. With this said, a number of factors have been identified that influence the effectiveness of gesture-based gaming. The first of these is the intuitiveness of the gestures used. In this regard, gestures should seem natural within the context of the game so that the user does not feel limited by their use~\citep{Kang04}. Another important factor is appropriate and immediate feedback so that the user feels that they have control over their environment. In addition, this allows users to calibrate the way in which they perform the various game-play gestures~\citep{Payn06}. This is related to the fact that gestural interfaces do not provide the same affordances as those present in traditional games~\citep{Kang04}. For this reason, the gestures in use should be simple and easy to learn~\citep{Payn06}. Finally, simple reminders or graphical hints should also be provided to aid recall and a consistent design should be used so that related commands share a gestural structure~\citep{Payn06}.

\citet{Payn06} found that previous gaming experience had a significant effect on how the gestural interface was received. Experienced gamers expected greater precision in interaction whereas inexperienced gamers felt that the gestural interface was easier to use and more reactive as it required less memorisation of key combinations~\citep{Payn06}. The differences in perception in this case may suggest that gesture-based gaming reduces the learning curve present in many games by making it easier for inexperienced players to become immersed in the environment~\citep{Payn06}.

Thus it appears that the appropriateness of gestures in gaming is dependant on the characteristics of the game and gestural interface themselves as well as the expectations of the player.

\subsection{Conclusion}
Gestures in games may serve as a novel interaction technique that is able to more rapidly immerse novices within a game-play environment. However, the cognitive demands surrounding learning and memorisation of gestures need to be carefully considered and affordances provided where possible. The choice of gestures should also be consciously examined within the context of the specific game being developed. Overall, however, further research is needed on this topic in order to fully understand the implications of the use of gestures in games.
 
